% 本文件由 LaTeX 排版 Agent 根据有效 translation.json 自动生成。
% author=Sulpitius Severus (c. 363-c. 420); work_id=3502; title=真确书信

\chapter{\WorkTitle{真确书信}}

% structure: p0001 source=h1 target=omitted reason=page-title-already-rendered-by-parent

% structure: p0002 source=h2 target=section reason=relative-heading-rank; base=h2

\section{第一封信。致\Person{欧瑟伯}。}

驳斥某些嫉妒的攻击\Person{玛尔定}的人。\par

\Emph{昨日}有几位隐修士来到我这里，在无尽的闲谈和冗长的对话中，提到了我就那位圣者\Person{玛尔定}的生活所发表的小作品，我很高兴听说许多人热切而仔细地阅读它。然而，与此同时，我得知某个人在邪灵的影响下，竟问说：玛尔定曾使死人复活、从火焰中拯救房屋，为何他自己最近却受制于火的力量，从而遭受危险的痛苦？说这话的人，无论他是谁，真是个可悲的人！我们从他口中认出了古时犹太人的恶语，他们曾这样辱骂挂在十字架上的主：\BibleQuote{他救了别人，却救不了自己。}（\BibleRef{玛 27:42}）的确，显然无论所指的人是谁，如果他生活在那个时代，他也会准备好用这些话攻击主，正如他以类似的方式亵渎主的圣者一样。那么，我问你，无论你是谁，情况如何？难道玛尔定真的没有能力、不分享圣洁，因为他曾暴露于火的危险吗？哦，你有福的人，你在一切事上都像宗徒们，甚至在你所承受的辱骂上也是如此！确实，据说外邦人对\Person{保禄}也曾抱有同样的想法，当毒蛇咬了他时，他们说：\BibleQuote{这个人必定是凶手，虽然从海里获救，命运却不许他活着。}（\BibleRef{宗 28:4}）但他把毒蛇抖到火里，没有受害。他们却以为他会突然倒下，很快死去；但见他没有受害，就改变了看法，说他是神。但是，你这最可怜的人，你本应从那例子中自己认识到你的虚假；因此，如果玛尔定似乎被火焰触碰而成为你的绊脚石，你反倒应该把他的仅仅被触碰归因于他的功德和能力，因为他在火焰包围中并没有灭亡。因为，你要承认，你这可怜的人，承认你似乎不知道的：几乎所有圣人在危险中所表现出的，比他们展示的德行更为卓越。我确实看到\Person{伯多禄}信德坚固，违反本性在海浪上行走，用脚印踏压不稳定的水。但外邦人的传道者（\Person{保禄}）被海浪吞没，三天三夜后，水又将他从深处托起，在我看来并不比他渺小。不，我几乎倾向于认为，在深海中存活比在海深处行走更为伟大。但是，你这愚昧的人，我想你没有读过这些事，或者读了却不明白。因为蒙福的圣史不会在圣经中记录那种事——除非出于天主的引导——（除非从这些事例中，人的心灵可以受到关于海难和蛇的危险的教导！）正如宗徒所叙述的，他以自己的赤身、饥渴和盗贼的危险为荣，这一切确实都是圣人们共同忍受的，但义人的卓越之处始终在于忍受并战胜这些事，在所有的考验中，他们忍耐且不可战胜，负担越重，他们就越英勇地克服它们。因此，这件归因于玛尔定软弱的事，实际上充满了尊严和荣耀，因为他经受了最危险的灾难，却以胜利者的姿态出现。但不要有人奇怪，在我写的关于他生活的那篇论文中，我遗漏了所述的事件，因为在那部作品中，我公开承认我没有包含他所有的行为；理由很充分：如果我打算全部叙述，我必定会向读者呈现一部庞大的著作。事实上，他的成就数量有限，无法全部写在一本书里。尽管如此，我不会让这个引起疑问的事件继续埋没，而将如实叙述整个事情，免得我似乎有意忽略那个可能被用来诽谤圣者的事。\par

时值仲冬，\Person{玛尔定}按主教巡视教区各堂的常例，来到某堂区。司铎们已在圣堂内僻静处为他准备了住所，并在朽坏且极薄的地板下生起大火。他们还用大量稻草为他铺了一张床榻。\Person{玛尔定}就寝时，因这过分柔软的卧榻而感不适——他素来只以一块粗麻布铺地而卧。于是，他仿佛受了亏待一般，将全部稻草掀到一旁。恰巧，他掀开的稻草有一部分落在火炉上。而他自己则照常疲乏地躺卧在光地上，因长途跋涉而疲惫不堪。约至夜半，火焰从本已不牢固的火炉中蹿起，引燃了干草。\Person{玛尔定}被这意外惊醒，迫于眼前的危险，但主要如他后来所述，因魔鬼的诡计与催逼，未能及时求助于祈祷。他想要逃出，却与闩住门的门闩长久费力地搏斗。不久他便察觉自己已被可怖的烈火包围，火焰甚至已引燃了他所穿的衣服。他终于恢复了一贯的信念——他的安全不在于逃跑，而在于\OriginalTerm{主}{Lord}，他拿起信德与祈祷的盾牌，将自己完全交托于主，卧于火焰之中。那时，火焰因天主的干预而退却，他继续在无害的火焰环中祈祷。守候在门外的隐修士们听到噼啪作响、挣扎的火焰声，砸开了闩着的门；扑灭火焰后，他们从火中救出\Person{玛尔定}，一直以为他在如此长久的火焰中早已烧成灰烬。如今主为我作证，他亲自向我讲述，并叹息着承认，他在这事上受了魔鬼诡计的欺骗，因为醒来时没有明智地用信德和祈祷抵御危险。他还补充说，当他心神不安地竭力开门时，火焰在他四周猛烧；而他一旦重新寻求十字架的援助，运用祈祷的武器，中间的火焰便退却，他感到火焰给他洒下甘露般的清凉，而此前他刚尝尽烈火烧身的痛苦。凡读到这封信的人，都当明白：\Person{玛尔定}确实经历了那危险，却真正蒙主悦纳地安然度过。\par

% structure: p0006 source=h2 target=section reason=relative-heading-rank; base=h2

\section{第二封信。致执事\Person{奥勒留}。}

\Person{苏尔皮提乌斯}看见圣\Person{玛尔定}的异象。\par

\Person{苏尔皮提乌斯·塞维鲁}向执事\Person{奥勒留}问安——\par

早晨你离开我之后，我独自坐在我的小室里；我心中常常涌起对未来的希望，以及对现世的厌倦、对审判的恐惧、对惩罚的畏惧，而这一切思想的源头，是我对我罪的记忆，这使我疲惫而可怜。于是，在因心灵的痛苦而疲惫不堪地躺到床上后，睡意悄悄袭来，这常常因忧郁而起；这样的睡眠，就像清晨时分总是轻浅而不确定，它只是以游移而朦胧的方式浸透我的肢体。因此，这不同于其他种类的睡眠，它使人能在几乎清醒时感到自己在做梦。在这样的情形下，我似乎突然看见\Person{圣玛尔定}以主教的身份向我显现，身穿白袍，面容如火，眼睛如星，头发呈紫色。他如此显现，带着我所熟悉的容貌和身形，以至于我几乎难以表达我的意思——他不能被定睛注视，却能被清楚认出。嗯，他向我投来温柔的微笑，右手拿着我写的那部关于他生平的小论著。我则抱住他神圣的膝盖，按惯例恳求他的祝福。于是，我感到他的手以最温柔的触摸放在我的头上，在庄严的祝福话语中，他一遍又一遍地重复那与他嘴唇如此熟悉的十字圣号。不久，当我热切地注视着他，无法满足于凝望他的面容时，他突然被从我这里带走，升向高处。最后，当他穿过广阔的空气时，我努力的目光追随着他在疾驰的云中上升，我再也看不见他。不久之后，我看见圣司铎\Person{克拉鲁斯}，\Person{玛尔定}的门徒，近日去世的，以我见他老师同样的方式升了上去。我厚颜无耻地想要跟随，当我瞄准并努力追求如此崇高的脚步时，突然醒来；从睡梦中惊醒，我开始为这异象而欢喜，这时一个家中的僮仆进来，面容比通常用言语表达悲伤的人更加忧郁。\Quote{什么，}我问他，\Quote{你要以如此忧郁的神情告诉我什么？}他回答说：\Quote{两个隐修士，刚从\Place{都尔}来，他们带来了\Person{玛尔定}去世的消息。}我承认我心碎了；我泪流满面，痛哭不已。不仅如此，即使现在，当我向你写下这些事时，兄弟，我的眼泪仍在流淌，对于这几乎无法承受的悲伤，我找不到任何安慰。我希望你，当这消息传到你那里时，能分担我的悲伤，就像你曾与我一同分享对他的爱一样。那么，我恳求你，立即到我这里来，让我们共同哀悼我们共同爱着的人。然而我非常清楚，这样的人不应被哀悼，他在战胜并得胜世界之后，终于获得了正义的冠冕。尽管如此，我无法控制自己而不悲伤。毫无疑问，我已先派了一位在天上为我辩护的人，但同时我也失去了此生中极大的安慰之源；然而，如果悲伤能让位于理性的影响，我确实应当喜乐。因为他如今已与宗徒和先知们为伍，（请容我对天上的圣人们表示敬意）在那义人的集会中，他不逊于任何人，正如我坚定地希望、相信并确信的，他特别加入了那些在羔羊的血中洗净自己衣裳的人的行列。他现在跟随羔羊作为他的向导，不受任何污点的玷污。因为尽管我们时代的特点未能确保他获得殉道的荣誉，但他绝不会缺乏殉道者的荣耀，因为无论从誓愿还是德行上，他都同样能够并愿意成为殉道者。但如果他曾在\Person{尼禄}和\Person{德西乌斯}的时代被允许参加当时进行的斗争，我以天地的天主作证，他会自由地屈服于酷刑架，并欣然将自己投身于火焰：是的，堪与那些杰出的希伯来青年相比，在环绕的火焰中，即使在火窑的中心，他也会唱一首赞美主的诗歌。但如果迫害者碰巧喜欢对他施加\Person{依撒意亚}所受的惩罚，他绝不会显得不如那位先知，也不会因被锯子和刀剑撕裂肢体而退缩。如果邪恶的狂怒更喜欢将这有福的人驱赶过陡峭的岩石或险峻的高山，我断言，他会紧握真理的见证，甘愿坠落。但如果，如同外邦人的导师的榜样（确实常常发生），他被列入其他被判刀剑之死的牺牲者中，他会最先催促刽子手执行工作，以便获得血冠。事实上，他远未逃避对主的宣认，面对所有那些常使人类软弱不堪的刑罚和惩罚，他会如此坚定不移，以至于在所受的苦难和折磨中带着喜乐和欢愉微笑，无论施加给他怎样的酷刑。但尽管他实际上并未遭受这些，他却完全没有流血就充分达到了殉道的荣誉。因为为了永生的希望，他忍受了何等人类痛苦的煎熬啊——饥饿、守夜、赤身、禁食、恶人的辱骂、恶徒的迫害、对弱者的关怀、对遇险者的焦虑？因为谁曾受苦而\Person{玛尔定}不与他同受呢？谁曾跌倒而他不燃烧呢？谁曾灭亡而他不深自哀恸呢？除了他每天与人类和属灵的邪恶进行各种争斗之外，当他总是受到各种试探攻击时，在他身上占上风的是征服中的坚韧、等待中的忍耐和承受中的平和。啊，人哪，在虔敬、慈悲、爱德上真是不可言喻，这些连在圣人中也因世界的冷漠而日渐冷淡，但在他身上却一直增长到底，日复一日地持久！我本人有幸在他内以卓越的程度享有了这恩宠，因为他以一种特殊的方式爱我，尽管我远不配得这样的情谊。而每当想起，我的眼泪再次迸发，呻吟从心底涌出。将来我在谁身上能找到像在他那里一样的灵魂安息呢？我又能在谁的爱中获得同样的安慰呢？我这个可怜的人，深陷苦难，如果我还有命活着，我能停止哀叹我比\Person{玛尔定}活得久吗？将来我生命中还会有任何快乐，或任何无泪的日子或时辰吗？或者，我最亲爱的兄弟，我能够向你提及他而不哀恸吗？然而，与你交谈时，我难道能谈论任何其他话题吗？但我为何激起你的眼泪和哀恸呢？所以我现在希望你得到安慰，尽管我无法安慰自己。他不会离开我们；相信我，他永远不会，永远不会抛弃我们，而是会在我们谈论他时与我们同在，在我们祈祷时亲近我们；甚至今天他屈尊赐予的幸福，即看见他在荣耀中，将来他会经常赐予；并且他会像刚才一样，以他不间断的祝福保护我们。然后，按照异象的安排，他表明了天堂是为跟随他的人敞开的，并教导我们应当跟随他到何处；他指示我们应当将希望指向什么对象，我们的心应当转向什么成就。然而，我的兄弟，该怎么办呢？因为，我自己很清楚，我永远无法攀登那艰难的斜坡，进入那些有福的区域。一个悲惨的重担如此沉重地压着我；当我因笼罩我的罪的重负而无法确保升天时，残酷的压力反而使我在苦难中沉沦于绝望之地。尽管如此，希望仍在，一个最后且唯一的希望：我靠自己无法获得的，也许能因\Person{玛尔定}为我所作的祈祷而被认为相称。但是，兄弟，为何我还要用一封如此絮叨的信占用你的时间，从而拖延你到我这里来呢？同时，我的篇幅已满，不能再容纳更多了。然而，我之所以将谈话延长到有些过分的程度，是为了这个目的：既然这封信向你传达了悲伤的消息，它也可能通过我与你的这种友好交谈为你提供安慰。\par

% structure: p0010 source=h2 target=section reason=relative-heading-rank; base=h2

\section{第三封信。致岳母\Person{巴苏拉}。}

\Person{圣玛尔定}如何从此世进入永生。\par

\Person{苏尔皮基乌斯·塞维鲁斯}致其敬爱的尊长\Person{巴苏拉}，问安。\par

如果子女可以合法地将父母传唤到法庭，那么我显然可以以抢劫和掠夺的罪名，用正义的绳缚将你拖到裁判官的审判台前。因为我为何不应抱怨你所加于我的损害？你没有留下任何一点手稿、任何一本书、甚至一封信在家——你竟如此洗劫所有类似之物，并将它们公之于世。如果我以亲切的笔调给朋友写些东西；如果我一边自娱一边口述些什么，同时又希望它保密，所有这些似乎在你尚未写就或说出之前就已到达你手中。显然，我的秘书们亏欠于你，因为我所写的任何琐碎之作都通过他们而让你知晓。然而，我不能对他们动怒，如果他们真的服从你，并因你的慷慨对他们施以特别影响而侵犯了我的权利，还始终记得他们属于你而非属于我。是的，唯有你是有罪之人——唯有你应当受到责备——因为你既设圈套陷害我，又用诡计诱骗他们，以至于不加筛选地，那些亲切之作或草率放出的文字，未经润色和修饰就送到你手中。因为，暂且不提其他著作，我敢问，我最近写给执事\Person{奥勒留}的那封信，如何能如此迅速地到达你手中？因为，当时我身处\Place{图卢兹}，而你住在\Place{特里尔}，因挂虑你儿子而远离故土，我想知道，你利用了何种机会得到了那封亲切的信？因为我已收到你的信，你在信中说，我应在同一封信中提及我们师长\Person{玛尔定}之死时，描述那位圣者离世的方式。仿佛我当初发出那封信是为了让他人阅读，而不仅是他本人；又仿佛我注定要承担如此重大的工作，以至于所有关于\Person{玛尔定}的应知之事，都特别通过我这个作者公之于众。因此，若你想知道关于那位圣主教临终的任何事，你应向当时在场的人询问。我本人已决定不给你写任何东西，以免你将我到处公开。不过，若你承诺不把我寄给你的内容读给任何人听，我愿用寥寥数语满足你的愿望。因此，我将把我所知悉的以下细节告知于你。\par

我必须说明，\Person{玛尔定}早在他去世之前就已预知自己离世的时刻，并告知弟兄们他的离世已在眼前。在此期间，发生了一件事，促使他前往\Place{龚达特}教堂。因为那间教堂的司铎们彼此不和，\Person{玛尔定}希望恢复和平，虽然他深知自己时日无多，但为了这个目的，他并不畏惧踏上旅程。事实上，他认为若能在他身后留下一个恢复和平的教会，那将是为他美德加冕的上好冠冕。于是，他带着那群常伴随他的众多而神圣的门徒出发了。他在一条河里看到许多水鸟正忙于捕鱼，并注意到贪婪的食欲正驱使它们频繁地攫取猎物。\Quote{这，}他喊道，\Quote{正是魔鬼行径的写照：它们埋伏等待粗心大意的人，在他们尚未知晓时捕捉他们；一旦得手，便吞噬他们的猎物，而且对吞噬之物永不满足。}然后，\Person{玛尔定}以言语中的奇迹之力，命令那些鸟离开它们游水的池塘，前往干旱荒凉之地；他用以驱逐魔鬼的同一权柄，如今对待这些鸟。于是，所有那些鸟聚集成群，离开河流，飞向山岭和森林，令许多看到\Person{玛尔定}拥有如此权柄——甚至能号令飞鸟——的人惊叹不已。他在所去的那个村庄或教堂停留了一段时间，当司铎们之间恢复和平后，他正考虑返回隐修院时，突然开始体力衰微，于是他召集弟兄们，告诉他们他即将离世。那时，哀伤和悲痛笼罩了所有人，他们齐声哀叹说：\Quote{亲爱的父亲，你为何要离开我们？或者，你能将我们托付给谁来应付我们的孤苦？凶猛的豺狼会迅速攻击你的羊群，当牧人被击倒后，谁能从它们的撕咬中拯救我们？我们确实知道，你渴望与基督同在；但你在天上的赏报是安全的，不会因延迟而减少；反而要怜悯我们这些你将遗弃的孤苦之人。}然后，\Person{玛尔定}被这些哀叹所触动——他始终确实充满怜悯——据说忍不住流泪；他转向主，只以以下话语回答那些围着他哭泣的人：\Quote{主啊，如果我对你的子民还有必要，我不畏惧劳苦：愿你的旨意承行。}就这样，他在渴望与爱之间徘徊，几乎不知自己更愿选择什么；因为他既不愿离开我们，也不愿与基督分离。然而，他不看重自己的意愿，也不将任何事留给自己，而是将自己完全交付于主的旨意和权能。你们难道不觉得你们听到他在我重复的以下寥寥数语中所说的话吗？\Quote{主啊，肉身的争战实在可怕，的确，我至今已战斗了足够久；但是，如果你命我继续为保护你的羊群而承受同样的劳苦，我不拒绝，也不以年迈推脱此任命。我完全属于你，我将完成你指派的一切任务，并只要你还命令，我就继续在你的旗帜下服役。是的，虽然对一位历经长久劳苦的老人而言，解脱是甜蜜的，但我的心智战胜了我的年岁，我不愿向年老屈服。但如果你此刻怜悯我的高龄，主啊，你的旨意对我而言是好的；你必将守护那些我为之担忧的人的安全。}啊，无人能言语描述的人，既不被劳苦征服，甚至不被死亡征服，他对两者都不偏不倚，既不畏惧死亡，也不拒绝生活！因此，尽管他连续几天发高烧，他仍然没有放弃天主的工作。他整夜祈祷守夜，迫使自己疲惫的肢体服侍他的精神，他躺在荣耀的床铺上，铺着麻布和灰烬。当他的门徒恳求他至少允许在身下铺些普通稻草时，他回答说：\Quote{基督徒死在灰烬中才合适；如果我留给你们不同的榜样，我就有罪了。}然而，他双手和双眼坚定地朝向天空，从未停止用祈祷释放他不屈不挠的精神。当围在他身旁的司铎们请求他稍微换换姿势以舒缓身体时，他喊道：\Quote{亲爱的兄弟，请允许我专注地仰望天堂，而不是大地，这样我即将踏上旅程的灵魂就能被引领向主。}说完这些话，他看见魔鬼站在近旁，喊道：\Quote{你这血腥的怪物，为什么站在这里？你在我身上找不到什么，你这致命的家伙：亚巴郎的怀抱就要接纳我了。}\par

他说完这些话，灵魂便离去了；当时在场的人向我们作证，他们看见他的面容好像天使的面容。他的四肢也洁白如雪，以致人们惊叹：\Quote{谁能相信这人曾身穿苦衣，谁又能想象他曾满身蒙灰？}因为此时他显出的容貌，仿佛已预现未来复活的光荣，并带着已经改变的身体本性。然而，几乎难以置信的是，参加他葬礼的人群是何等众多：全城的人都涌出来迎接他的遗体；整个地区和村庄的居民，还有许多来自邻近城市的人，都来送葬。啊，众人的悲痛何其大！哀悼的隐修士们尤其是多么深沉地哀哭！据说那一天聚集了几乎两千名隐修士——这是\Person{马丁}特有的光荣：因着他的榜样，如此众多的苗裔为事奉主而生长出来。毫无疑问，那时牧者正带领他自己的羊群前行——那群苍白的圣洁群众——身披斗篷的队伍，其中既有生命劳苦已结束的长者，也有刚刚向基督宣誓从军的年轻战士。还有贞女的歌咏团，因端庄而克制哭泣；她们以何等神圣的喜乐隐藏了自己的痛苦！无疑，信德阻止了流泪，但亲情却逼出了叹息。因为对他所得光荣的欢欣与对他死亡的虔敬悲伤并存。人们会原谅哭泣的人，也会祝贺喜乐的人，而每个人都宁愿自己悲伤，却愿他人喜乐。就这样，这群唱着天上赞歌的群众，护送着圣人的遗体直到墓地。让这景象与——我不说世俗的葬礼，甚至与凯旋——相比，有什么能与马丁的葬礼相提并论呢？让你们的世俗伟人将俘虏双手反绑在战车前示众吧。那些在马丁引领下战胜了世界的人，护送着马丁的遗体。让疯狂以万民的齐声赞美来荣耀这些尘世战士吧。马丁却以神圣的圣咏受到赞美，马丁以天上的赞歌受到尊荣。那些世俗之人在此世的凯旋之后，将被投入残酷的地狱深处，而马丁则被欢欣地接入亚巴郎的怀抱。马丁在地上贫穷卑微，却获准进入天堂，成为富足的进入。我相信，从那个福地，他作为我的守护者，在我写下这些事的时候注视着我，也注视着你们这些读者。\par

\SourceRecord{Translated by Alexander Roberts.；Nicene and Post-Nicene Fathers, Second Series；Vol. 11.；Edited by Philip Schaff and Henry Wace.；Buffalo, NY: Christian Literature Publishing Co.,；1894.；<http://www.newadvent.org/fathers/3502.htm>.}
